\documentclass[11pt]{article}

\usepackage{amsmath}
\usepackage{amssymb}
\usepackage{booktabs}
\usepackage{luatexja}

\title{輸出プラットフォーム型 vs 市場追求型のアイデア}
\date{}

\begin{document}

\maketitle

私たちが考えているのは、企業によるEU域内投資の志向性だと思います。この場合、売上データがないという問題をどう扱うかを考える必要があります。各投資について企業の意図を示す東洋経済データを利用できるのであれば別ですが、そのデータがどの程度信頼できるのか、あるいは他の研究者が使っているのかは分かりません。そのため、いくつかの代理変数を作ることを試みる必要があります。私／ChatGPT にはいくつかアイデアがあります。

\section{生産・販売分離指標}

最初の指標は、生産拠点と販売拠点の地理的な分離を捉えるものです。

以下のように定義します。

\[
M_{ij} = \text{企業 } i \text{ が国 } j \text{ に保有する製造現地法人の数}
\]

\[
S_{ij} = \text{企業 } i \text{ が国 } j \text{ に保有する販売・卸売現地法人の数}
\]

輸出プラットフォーム指標を次のように定義します。

\[
EP_i =
\frac{\text{販売現地法人はあるが製造現地法人はない国の数}}
{\text{現地法人を保有する国の総数}}
\]

この指標の値が高い場合、企業が限られた国で生産を行い、多くの市場に製品を販売していることを示すため、輸出プラットフォーム型の行動を示唆します。

\section{生産集中度と市場カバレッジの指標}

第2の方法は、製造活動と販売活動の地理的な集中度を比較するものです。

製造活動の集中度は、ハーフィンダール指数を用いて測定できます。

\[
H_M = \sum_j
\left(
\frac{M_{ij}}{\sum_j M_{ij}}
\right)^2
\]

同様に、販売活動の集中度は次のように定義できます。

\[
H_S = \sum_j
\left(
\frac{S_{ij}}{\sum_j S_{ij}}
\right)^2
\]

輸出プラットフォームのシグナルは、次のように定義されます。

\[
EP_i = H_M - H_S
\]

製造活動が販売活動よりもかなり集中している場合（\(H_M > H_S\)）、このパターンは輸出プラットフォーム型FDIと整合的です。

\section{生産ハブ指標}

もう1つ有用な指標は、多国籍企業ネットワーク内の生産ハブを識別するものです。企業が現地販売法人を持たない国に製造現地法人を保有している場合、その工場は他の欧州市場に輸出している可能性が高いと考えられます。

次のように定義します。

\[
ExportPlatformShare_i =
\frac{\text{現地販売法人を伴わない製造現地法人の数}}
{\text{製造現地法人の総数}}
\]

この値が高い場合、生産施設が現地ホスト国市場ではなく、外部市場にサービスを提供している可能性が高いことを示します。

\section{構造例}

表 \ref{tab:example} は、典型的な輸出プラットフォーム型の構造を示しています。

\begin{table}[h]
\centering
\begin{tabular}{lcc}
\toprule
国 & 製造現地法人 & 販売現地法人 \\
\midrule
アイルランド & あり & あり \\
ドイツ & なし & あり \\
フランス & なし & あり \\
イタリア & なし & あり \\
スペイン & なし & あり \\
\bottomrule
\end{tabular}
\caption{輸出プラットフォーム型FDIの構造例}
\label{tab:example}
\end{table}

この例では、アイルランドが生産ハブとして機能し、他国の販売現地法人が欧州市場全体に財を流通させています。

\section{なぜ東洋経済データでこの方法が使えるのか}

東洋経済のようなデータセットは、現地法人レベルで所在地、産業分類、親会社の識別情報を提供しています。輸出フローが観察できない場合でも、生産現地法人と販売現地法人の地理的な分離を利用することで、研究者は輸出プラットフォーム戦略を推測できます。

\section{実証研究における典型的な説明}

実証研究では、典型的には次のように説明できます。

\begin{quote}
現地法人レベルの輸出データが利用できないため、輸出プラットフォーム戦略は、生産現地法人と販売現地法人の地理的な分離を用いて代理する。製造活動が限られた国に集中している一方で、複数の欧州市場に販売子会社を維持している企業は、輸出プラットフォーム型FDIを行っている企業として分類される。
\end{quote}

以下の文献を見るとよいかもしれません。

https://www.sciencedirect.com/science/article/abs/pii/S0969593105000971

「本研究は、貿易ブロック内の小国経済における外国所有子会社の役割と特徴を検討している。外国所有子会社の役割として、ローカル・サテライト、切り詰められたミニチュア複製型、輸出プラットフォーム型、地域または世界レベルで権限を与えられたハブ型の4類型が識別されている。オランダの電子・電気応用産業における外国所有子会社の販売先の焦点について証拠が示されている。分析の結果、外国所有子会社のほぼ半数が、オランダだけでなく、より多くの市場に販売の焦点を置いていることが分かった。これらの子会社はオランダを輸出拠点として利用している。複数の生産工場を用いるのではなく、1つの拠点で生産された製品を欧州市場に供給している。結果は、非欧州の親会社を持つこれらの子会社が比較的若く、製造活動にも重点を置いていることを示している。」

\end{document}
